%\documentclass[a4paper, 10pt, twocolumn, journal, twoside]{IEEEtran}
\documentclass[a4paper, 12pt, onecolumn, journal, twoside, draftclsnofoot]{IEEEtran}

\newcommand*{\DOUBLECOL}{}	% turns this on for double column
\newlength{\subfigureheight}
\newlength{\figurewidth}
\ifdefined\DOUBLECOL
\setlength{\subfigureheight}{2.45in}
\setlength{\figurewidth}{3.5in}
\else
\setlength{\subfigureheight}{2.35in}
\setlength{\figurewidth}{6.5in}
\fi
\usepackage{amsfonts,amsmath,amssymb,bm,cancel,color,comment,epsfig,enumerate,empheq,floatrow,hyperref,mathdots,mathrsfs}
\usepackage{placeins}%\usepackage{algorithmic}
\usepackage[vlined,ruled,linesnumbered]{algorithm2e}
\usepackage{graphicx}
%\usepackage{amsthm}
\usepackage[super]{nth}
\usepackage[noadjust]{cite}
\usepackage{subcaption}
\usepackage{tabularx}
\usepackage{tabularray}		% make thing horizontal lines in tabular
\UseTblrLibrary{booktabs}	% make thing horizontal lines in tabular

\graphicspath{{eps/}}

\hypersetup{
	colorlinks=true,
	linkcolor=red,
	citecolor=green,
	filecolor=magenta,
	urlcolor=blue
}

\newenvironment{Ualgorithm}[1][htpb]{\def\@algocf@post@ruled{\kern\interspacealgoruled\hrule  height\algoheightrule\kern3pt\relax}%
\def\@algocf@capt@ruled{under}
\begin{algorithm}[#1]}
{\end{algorithm}}

\setcounter{MaxMatrixCols}{30}
\DeclareSymbolFont{grb}{OML}{cmm}{b}{it}
\DeclareMathSymbol{\zetab}{\mathord}{grb}{"10}
\DeclareMathSymbol{\etab}{\mathord}{grb}{"11}
\DeclareMathSymbol{\thetab}{\mathord}{grb}{"12}
\DeclareMathSymbol{\kappab}{\mathord}{grb}{"14}
\DeclareMathSymbol{\lambdab}{\mathord}{grb}{"15}
\DeclareMathSymbol{\mub}{\mathord}{grb}{"16}
\DeclareMathSymbol{\nub}{\mathord}{grb}{"17}
\DeclareMathSymbol{\rhob}{\mathord}{grb}{"1A}
\DeclareMathSymbol{\sigmab}{\mathord}{grb}{"1B}
\DeclareMathSymbol{\taub}{\mathord}{grb}{"1C}
\DeclareMathSymbol{\phib}{\mathord}{grb}{"1E}
\DeclareMathSymbol{\psib}{\mathord}{grb}{"20}
\DeclareMathSymbol{\omegab}{\mathord}{grb}{"21}
\DeclareMathSymbol{\epsilonb}{\mathord}{grb}{"22}
\DeclareMathSymbol{\varphib}{\mathord}{grb}{"27}
\input{Symbol_Shortcut.tex}

\SetKwRepeat{Do}{do}{while}%

\newtheorem{theorem}{\bf Theorem}		% by Victor
\newtheorem{corollary}{\bf Corollary}	% by Victor
\newtheorem{definition}{\bf Definition}	% by Victor

%\centerfigcaptionstrue

%\pagestyle{empty}

\makeatletter
\def\ifundefined{\@ifundefined}
\makeatother

\setlength{\parskip}{0ex}

% generating the double accent on o, i.e. \fH{o}, e.g. Erdos Renyi 
\newcommand\fH[1]{\sbox0{#1}\dimen0=\ht0 \advance\dimen0 -1ex
  \sbox2{\'{}}\sbox2{\raise\dimen0\box2}%
  {\ooalign{\hidewidth\kern.1em\copy2\kern-.5\wd2\box2\hidewidth\cr\box0\crcr}}}

\title{\vspace{-0.2em} \huge Latency Minimization for Uplink RSMA System}

\ifdefined \BLIND
\author{Carrson C. Fung $^1$ \\ %and Author 2$^2$ \\
$^1$ NYCU  \\
%email: \emph{Author 1's email, Author 2's email} \\
\thanks{Funding sources.}
}
\else
\author{Shanik Hubert, ~\IEEEmembership{Student Member,~IEEE}, and Carrson C. Fung$^*$,~\IEEEmembership{Member,~IEEE} \\
%Institute of Electronics, College of ECE, National Chiao Tung University, Hsinchu, Taiwan 300 %\\
%Voice: +886-3-573-1862$^2$ Fax: +886-3-572-4361$^2$ email: \emph{rusly.eic04g@nctu.edu.tw, c.fung@ieee.org \vspace{-1.5em}} 
\thanks{This work has been partially supported by Google exploreCSR Grant 111Q90004C.

Shanik Hubert and Carrson C. Fung  are both affiliated with the Department of Electrical and Computer Engineering, National Yang Ming Chiao Tung University, Hsinchu 300, Taiwan. (e-mail: dmytro.ivakhnenkov@gmail.com, c.fung@ieee.org)}
}
\fi


\markboth{IEEE Trans. on Communications}{Latency Minimization for Uplink RSMA System}
\begin{document}
\maketitle

%******************************************************************
\begin{abstract}
%******************************************************************
\end{abstract}

\begin{IEEEkeywords}
\end{IEEEkeywords}

%%%%%%%%%%%%%%%%%%%%%%%%%%%%%%%%%%%%%%%%%%%%%%%%%%%%%%%%%%%%%%%%%%
\section{Introduction}	\label{intro}
%%%%%%%%%%%%%%%%%%%%%%%%%%%%%%%%%%%%%%%%%%%%%%%%%%%%%%%%%%%%%%%%%%


  
%%%%%%%%%%%%%%%%%%%%%%%%%%%%%%%%%%%%%%%%%%%%%%%%%%%%%%%%%%%%%%%%%%
\section{Methodology} \label{sec_meth}
%%%%%%%%%%%%%%%%%%%%%%%%%%%%%%%%%%%%%%%%%%%%%%%%%%%%%%%%%%%%%%%%%%

\subsection{Uplink Blocklength Minimization for Reduction of Asynchronicity}
%******************************************************************
For MIMO uplink (UL), the receive signal is written as
\begin{equation*}
	\y = \sum_{k=1}^K  \H_k \F_k \s_k + \w \in \complex^{N_R},
\end{equation*}
where 
$\H_k \in \complex^{N_R \times N_k}$ denotes the uplink channel between the $k$th user and the AP, $\F_k \in \complex^{N_k \times d_k}$ denotes the precoding matrix for the $k$th user, $d_k$ denotes the number of spatial streams user $k$ is sending to the AP (and $d_k \leq N_k$), $\s_k \in \complex^{d_k}$ denotes the transmitted signal for the $k$th user, $\w \in \complex^{N_R}$  denotes the AWGN, with $N_k$ denoting the number of antennas at the $k$th user and $N_R$ denotes the number of antennas at the AP.  It is assumed that $E \left[ \s_k \s_k^H \right] = \I_{N_k}$.  Let 
\begin{equation*}
	x_k = \F_k \s_k,
\end{equation*}
so
\begin{align*}
	E \left[ \left| \x_k \right\|_2^2 \right] &= \F_k E \left[ \s_k \s_k^H \right] \F_k^H \\
	&= \F_k \F_k^H  \\
	&=  \left\| \F_k \right\|_F^2  \leq P_k,
\end{align*}
where $P_k$ is the transmit power from the $k$th user.  If SIC is used at the receiver, then the achievable rate at the receiver is
\begin{equation*}
	R_{sum} = \sum_{k=1}^K  \log_2 \det \left(  \I_{N_R}  +  \H_k \F_k     \left(  \sigma_w^2 \I_{N_R}  + \sum_{j \neq k}  \H_j \F_j \F_j^H \H_j^H  \right)^{-1}     \F_k^H \H_k^H  \right) 
\end{equation*}

Then it seems the latency minimization problem can be formulated as
\begin{equation}
	\begin{split}
		\min_{\F_k, B_k, \forall k}  &\;  \sum_{k=1}^K  \frac{ B_k } { R_{sum}  }   \text{ ( I think } R_k  \\
		s.t. &\;  \left\| \F_k \right\|_F^2  \leq P_k,  \forall k, \\
		&\; B_k \in \left\{  B_{k,\min}, \ldots, B_{k,\max} \right\},
	\end{split}
\label{eq_prob_form_basic}
\end{equation}
where $R_k$ is the Shannon rate for the $k$th user.  (\ref{eq_prob_form_basic}) makes sense when  infinite size blocklength is assumed.  However, this assumption is usually incorrect when latency is minimized as the design variable should be the number of blocks (i.e. channel used) for transmission, implying the blocklength is short and finite.  Possible extensions for (\ref{eq_prob_form_basic}) include optimizing not with the sum rate, but using rate such as
\begin{enumerate}[1.]
	\item Max-min rate: 
	\begin{equation*}
		\max_{\F_k, \forall k}  \min_k  R_k.
	\end{equation*}
This has a strong fairness guarantee, often leads to equal rates across all users if we can find a feasible solution.  This is the worst-case that Carrson was talking about.
	\item The proportional fairness:
	\begin{equation*}
		\max_{\F_k, \forall k}   \sum_{k=1}^K  \log (\R_k), 
	\end{equation*}
which balances the throughput and fairness. 
	\item Weighted fairness 
	\begin{equation*}
		\max_{\F_k, w_k, \forall k}  w_k \log (R_k), 
	\end{equation*}
where $w_k$ are weights for QoS differentiation, i.e., higher $w_k$ is, more resource will be allocated the $k$th user.
\end{enumerate}


\subsection{Uplink Latency Minimization Using RSMA}
%******************************************************************
Using RSMA for uplink transmission can theoretically acheive the optimal rate region according to \cite{rimoldi1996} by splitting the original $K$ input channels to $2K-1$ inputs (though the $K$th user message need not be split, but most papers consider splitting the last user message for simplicity).  The transmitted message of the $k$th user $\x_k$ is split into 
\begin{equation*}
	\x_k = \sum_{i=1}^2  \\F_{ki} \s_{ki},
\end{equation*}
where $i$ denotes the stream index.  So the received signal at the AP can be written as
\begin{equation*}
	\y = \sum_{k=1}^K \sum_{i=1}^2 \H_{ki} \F_{ki} \s_{ki} + \w \in \complex^{N_R},
\end{equation*}
where
\begin{equation*}
    \sum_{i=1}^2 \left\| \F_{ki} \right\|_F^2  \leq P_k
\end{equation*}

SIC is assumed to be used to decode each user's signal at the AP, which involves the use of a decoding order $\pi$, let the permutation $\pi$ belong to the set $\Pi$ defined as the set of all possible decoding order of all $2K$ messages from $K$ users.  Let the decoding order of message $s_{ki}$ from the $k$th user to be $\pi_{ki}$, then the achievable rate of decoded message $s_{ki}$ can be given as
\begin{equation*}
	\text{\textcolor{red}{missing equation here}.}
\end{equation*}
Let $\mcK$ be the set of user devices, Let $\mcJ = \{1, 2 \}$ be the set of containing the number of split streams.  Then the rate of the $i$th stream for user $k$ is
\begin{equation*}
	R_{ki} = \log_2 \det \left(  \I_{N_R}  +  \H_k \F_{ki}     \left(  \sigma_w^2 \I_{N_R}  + \sum_{\{(j{\in}\mathcal{K},i{\in}\mathcal{J})|\pi_{ji}{>}\pi_{ki}\}}  \H_j \F_{ji} \F_{ji}^H \H_j^H  \right)^{-1}     \F_{ki}^H \H_k^H  \right), 
\end{equation*}
where $\{(j{\in}\mathcal{K},i{\in}\mathcal{J})|\pi_{ji}{>}\pi_{ki}\}$ is the set that represents  the messages $s_{ji}$ that are decoded after  message $s_{ki}$.  Hence the total rate for user $k$ is
\begin{equation*}
    R_k = \sum^2_{i=1} R_{ki}
\end{equation*}
To optimize the rate in uplink RSMA, a joint optimization of both the power management and message decoding order must be done, which results in solving 
%Max-Min Rate
\begin{equation*}
	\begin{split}
		\max_{\F_k, \forall k}  \min_{k, \pi} &\; R_k.  \\
    	\mathrm{s.t.} &\; \sum_{i=1}^2 \left\| \F_{ki} \right\|_F^2  \leq P_k \\
    	&\; \pi \in \Pi, \left\| \F_{ki} \right\|_F^2 \geq  0, \forall k \in \mcK, j \in \mcJ 
    \end{split}
\end{equation*}
This is similar to the optimization problem modeled in \cite{yang2019} where they optimize for the sum rate.


\subsection{Uplink Latency Minimization Using Finite Blocklength Minimzation}
%******************************************************************
To consider latency in URRLC systems with finite blocklength, let us model channel rate not in terms of (bits/s) but as (bits/channel use).  Let blocklength of the $k$th user be $n_k$ or the number of channel used per symbol. Let $T_s$ be the time taken to transmit one symbol.  Then the latency of the $k$th user can be expressed as
\begin{equation*}
	\begin{split}
		\text{Latency}_k &= n_k \times T_s \\ 
		&= \frac{B (\text{bits})}{R\text{(bits/channel use)}} \times T_s (\text{sec})
\end{split}
\end{equation*}
Thus latency for the $k$th user is optimized by optimizing for $n_k$. 

According to \cite{durisi2016}, the \textbf{rate under finite blocklength} is expressed as
\begin{equation}
	R^*(n,\epsilon) = C - \sqrt{\frac{V}{n}}Q^{-1}(\epsilon) + \mathcal{O}\left(\frac{\log n}{n}\right), 
\label{eq_rate_finite}
\end{equation}
where $C$ is Shannon capacity of channel, $Q^{-1}(\cdot)$ is the inverse of the Gaussian $Q$ function, $V$ is the channel dispersion.  According to \cite{durisi2016}, (\ref{eq_rate_finite})  is used to sustain the desired error probability $\epsilon$ for a given packet size $n$, one incurs a penalty on the rate (compared to the channel capacity) that is proportional to $\frac{1}{\sqrt{n}}$.

According to \cite{you2023}, the \textbf{rate under finite blocklength in massive MIMO systems} in high SNR regime is approximated as 
\begin{equation*}
	R_k \approx d_k\log\left(1+\rho_k\right)-\sqrt{\frac{d_k}{n_k}}Q^{-1}\left(\varepsilon_k\right).
\end{equation*}
which is the rate of user $k$.  $\rho_k$ is the SNR of user $k$ and $Q^{-1}(\cdot)$ denotes the inverse of Gaussian $Q$ function.  At this point, we are not sure why \cite{you2023} removed the channel dispersion $V$, which is part of (\ref{eq_rate_finite}).  The inverse of the $Q$ function is  
\begin{equation*}
	Q^{-1}(x) \overset{\Delta}{\operatorname*{=}}\int_{x}^{\infty}\frac{1}{\sqrt{2\pi}}e^{-t^{2}/2} dt.
\end{equation*}
Rearranging the equation, when $d_k$ is large enough, 
\begin{equation*}
    n_k\approx\frac{d_k\left[Q^{-1}\left(\varepsilon_k\right)\right]^2}{\left[d_k\mathrm{log}_2\left(1+\rho_k\right)-R_k\right]^2}\propto\frac{1}{d_k},
\end{equation*}
This shows that the latency in the time domain can be reduced by increasing the DoFs in the space domain, for the purpose of maintaining a certain level of performance target. According to \cite{you2023} 
\begin{equation}
	\begin{split}
    	n_k &= B_k/R_k \\
		&= \frac{B_k}{d_k\log\left(1+\rho_k\right)-\sqrt{\frac{d_k}{n_k}}Q^{-1}\left(\varepsilon_k\right)}. 
	\end{split}
\label{eq_blklength_MIMO}
\end{equation}
Solving (\ref{eq_blklength_MIMO}) would give another equation that describe $n_k$ if we know the number of bits $B_k$. We am not sure how this fits in with the other formulas.  From (\ref{eq_blklength_MIMO}), it can also be seen there are several ways to decrease $n_k$, which are 
\begin{enumerate}
    \item Increase spatial DoF $d_k$
    \item Increase SNR $\rho_k$
    \item Increase error probability $\epsilon_k$ (reliability)
\end{enumerate}

The \textbf{min-max blocklength/Latency} can be expressed as
\begin{equation*}
    \min_{\forall k}  \max_{k,R_k,\epsilon_k, }  n_k. 
\end{equation*}

\subsection{Uplink Latency Minimization with RSMA and Finite Blocklength Minimization}
%******************************************************************
\cite{zhang2023} considered minimizing the blocklength in uplink RSMA in a simple 2-user system. However in our problem, we want to improve the latency of the worst-case user, i.e. the user with the worst channel (e.g. worst channel-to-noise (CNR)), we should consider increasing $d_k$, $\rho_k$, or $\epsilon_k$ of the weakest user through precoder design, while decreasing the same parameters of the stronger users so that all users will experience the same latency, thus reducing asynchronality among the users.   
%
%I feel this has a naturally synergestic approach, as eventhough you increase the error probability of the weaker user, the weighing mechanism in Fed-async takes care to reduce the impact of data from these users on the training.
%
Similarly reducing the error probability of the stronger user will be helpful as the data from teh stronger users are weighted higher. This helps to reduce the unfariness of asynchrononicity and in the context of federated learning, this ensures that the weaker users still have an impact on the training while not weighting the data higher due to higher error rate fro these users due to the poor channel.  This is something that i do not believe is explored in \cite{zhang2023}.

\subsection{Problem Formulation of Blocklength Reduction to Reduce Asynchronality}
%******************************************************************
\textbf{Written  on 8-21-2025.} \\

The Shannon rate of the receiver can be defined as,
\begin{equation*}
	R_{k} =  \log_2 \det \left(  \I_{N_R}  +  \H_k \F_k     \left(  \sigma_w^2 \I_{N_R}  + \sum_{j \neq k}  \H_j \F_j \F_j^H \H_j^H  \right)^{-1}     \F_k^H \H_k^H  \right) 
\end{equation*}
Defining the latency of the $k$th use as 
\begin{equation*}
    T_{k} = \frac{n_k}{f_s}, 
\end{equation*}
where $T_{k}$ is the latency, $n_k \in \mathbb{Z}_+$ is the number of channel used and $f_s$ is the symbol rate.  Considering error probability $\epsilon$ and using the it to calculate latency under the mean no. of retranmsmissions 
\begin{equation*}
    \begin{split}
        &\text{P}(n \ \text{transmissions}) = \epsilon_k^{n-1}(1-\epsilon_k). \\
        &\mathbb{E}( \text{no. of transmissions}) = n_{mean} = \frac{1}{1-\epsilon_k}.
    \end{split}
\end{equation*}
Sincer a finite blocklength $n_k$ is considered, the upper bound to the rate of system for the $k$th user is written as \cite{durisi2016}
\begin{equation*}
	\frac{B_k}{n_k} \leq R_k^*(n, \epsilon) = R_k - \sqrt{\frac{V_k}{n_k}}Q^{-1}(\epsilon_k), 
\end{equation*}
where $B_k$ is the fixed number of bits that need to sent by the $k$th user and $n_k$ is the number of blocklength or channel uses or symbols for the $k$th user. Thus $R_k^*(n, \epsilon)$ can be defined in the unit of (bits/channel use).  Then the weighted latency minimization problem can be formulated as
\begin{equation}
	\begin{split}
		\min_{\F_k, n_k, \epsilon_k} &\; \sum_{k \in K} W_k \frac{n_k}{f_s} \\
		\text{s.t.} &\; \frac{B_k}{n_k} \leq R_k(\F_k) - \sqrt{ \frac{V_k(\F_k)}{n_k} } Q^{-1} (\epsilon_k), \quad \forall k \in \mcK\\
		&\; n_k \in [n_{\text{min}}, n_{\text{max}}], \quad \forall k \in \mcK \\
		& \; \|\F_k\|_F^2 \leq P_k, \quad \forall k \in \mcK.
	\end{split}
\label{eq_short_blocklength_latency_min_prob}
\end{equation}
A more complete formulation by considering user specific $\epsilon_k$ and retransmission affected latency can be considered as
\begin{equation}
	\begin{split}
		\min_{\F_k, n_k, \epsilon_k} &\; \sum_{k \in K} W_k \frac{n_k}{f_s} \frac{1}{1-\epsilon_k} \\
		\text{s.t.} &\; \frac{B_k}{n_k} \le (1-\epsilon_k) \Bigg[ R_k(\F_k) - \sqrt{\frac{V_k(\F_k)}{n_k}} \, Q^{-1}(\epsilon_k) \Bigg], \quad \forall k \in \mcK, \\
		&\; n_k \in [n_{\min}, n_{\max}], \quad \forall k \in \mcK, \\
		&\; \epsilon_k \in [0, \epsilon_0], \quad \forall k \in \mcK, \\
		&\; \|\F_k\|_F^2 \le P_k, \quad \forall k \in \mcK.
	\end{split}
\label{eq_short_blocklength_latency_min__retrans_prob}
\end{equation}


\begin{comment}
\subsection{Applying QoS constraint to the user rate}
%******************************************************************
\textbf{Written on 8-21-2025}\\
This was the original equation which shows the relationship between the user rate $\frac{B_k}{n_k}$ and the theoretical upper limit $R^*(n_k, \epsilon_k) = R_k - \sqrt{\frac{V_k}{n_k}}Q^{-1}(\epsilon_k) $.
\begin{equation*}
    \frac{B_k}{n_k} \leq R_k - \sqrt{\frac{V_k}{n_k}}Q^{-1}(\epsilon_k) \\
\end{equation*}

We then added a QoS bound $\gamma_k$ to give a lower bound to $R^*(n_K,\epsilon_k)$ that is not a function of the variables $n_k$$ as 
\begin{equation*}
     \frac{B_k}{n_k} \leq R_k - \sqrt{\frac{V_k}{n_k}}Q^{-1}(\epsilon_k) \\ \geq \gamma_k
\end{equation*}

This introduces 2 lower bounds for $R^*(n_k \epsilon_k) \geq$  $\frac{B_k}{n_k}$ and $R^*(n_k \_k) \geq$  $\gamma_k$. We know that $\frac{B_k}{n_k}$ is monotonically decreasing in $n_k$, while $R^*(n_k,\epsilon_k)$ is monotonically increasing in $n_k$.  Since the objective is to decrease $n_k$, find an optimal solution such that the equality condition between $\frac{B_k}{n_k}$ and  $R^*(_k)$ is true, i.e.
\begin{equation*}
    \frac{B_k}{n^*_k} = R^*(n^*_k,\epsilon_k) = R_k - \sqrt{\frac{V_k}{n^*_k}}Q^{-1}(\epsilon_k)
\end{equation*}
%let  increase until it meets the decreasing in equality. Now assuming that for the optimal value $n_k^*$ (this is the smallest possible value of $n$ we can get, ie: min latency), 
Then there are 2 cases that must be considered
\begin{enumerate}
    \item Case I: $\gamma_k > R^*(n^*_k,\epsilon_k)$ \\
    In this case, since $\gamma_k > R^*(n^*_k,\epsilon_k)$, $R^*(n^*_k,\epsilon_k)$ is never achieved. \\
    
    As $n_k$ keeps decreasing, $R^*(n_k, \epsilon_k)$ decreases until,
    \begin{equation*}
        R^*(n_k,\epsilon_k) = R_k - \sqrt{\frac{V_k}{n_k}}Q^{-1}(\epsilon_k) = \gamma_k
    \end{equation*}
    Solving for $n_k$
    \begin{equation*}
    n_k = V_k\left(\frac{ Q^{-1}(\epsilon_k)}{R_k - \gamma_k}\right)^2
    \end{equation*}

    This is the minimum blocklength achievable, thus in this case, 
    $n^*_k = V_k \left(\frac{Q^{-1}(\epsilon_k)}{R_k - \gamma_k}\right)^2$
    and maximum user rate = $\frac{B_k}{n^*_k}$

    \item Case II: $\gamma_k < R^*(n^*_k, \epsilon_K) = \frac{B_k}{n^*_k}$ \\
    In this case, $\frac{B_k}{n^*_k}$ provides a lower bound to $R^*(n_k,\epsilon_k)$ before $\gamma_k$. \\

    Thus we believe the equation can be rewritten with $\gamma_k$ acting as a lower bound for $\frac{B_k}{n_k}$ and $R^*(n_k, \epsilon_k)$,
	\begin{equation*}
    	\gamma_k \leq \frac{B_k}{n_k} \leq R^*(n_k, \epsilon_k) = R_k - \sqrt{\frac{V_k}{n_k}}Q^{-1}(\epsilon_k)
	\end{equation*}

    We must solve this to obtain $n^*_k$.
\end{enumerate} 

Solution for (\ref{eq_short_blocklength_latency_min_prob}) can be obtained via alternating optimization.  

We can solve this by first evaluating, 

\begin{equation*}
\begin{split}
    R_k - \sqrt{\frac{V_k}{n_k}}Q^{-1}(\epsilon_k) &= \gamma_k\\
    n^*_k &= V_k \left(\frac{Q^{-1}(\epsilon_k)}{R_k - \gamma_k}\right)^2
\end{split}
\end{equation*}

This is not our final $n^*_k$, once we obtain this $n^*_k$, we substitute it in $\frac{B_k}{n_k}$, \\
\\
if:
\begin{equation*}
    \frac{B_k}{n^*_k} < R_k - \sqrt{\frac{V_k}{n^*_k}}Q^{-1}(\epsilon_k)
\end{equation*}

 then: According to Case I, \\
\begin{equation*}
n^*_k = V_k \left(\frac{Q^{-1}(\epsilon_k)}{R_k - \gamma_k}\right)^2
\end{equation*}
 \\else if: 
\begin{equation*}
    \frac{B_k}{n^*_k}> R_k - \sqrt{\frac{V_k}{n^*_k}}Q^{-1}(\epsilon_k)
\end{equation*}
 
 then: According to Case II, 
 
the inequality of $\frac{B_K}{n_K} \leq R^*(n_k,\epsilon_k)$ is violated, thus to obtain true $n^*_k$ we must solve for,
\begin{equation*}
    \gamma_k \leq \frac{B_k}{n_k} \leq R_k - \sqrt{\frac{V_k}{n_k}}Q^{-1}(\epsilon_k)
\end{equation*}

\textbf{Problem Formulation with QoS constraint}
\begin{equation*}
	\begin{split}
		\min_{F_k, n_k, \gamma_k, \epsilon_k} &\; \sum_{k \in K} W_k \frac{n_k}{f_s} \\
		\text{s.t. } &\; \frac{B_k}{n_k} \le   R_k - \sqrt{\frac{V_k}{n_k}} \, Q^{-1}(\epsilon_k)  \geq \gamma_k, \quad \forall k \in K, \\
		&\; n_k \in [n_{\min}, n_{\max}], \quad \forall k \in K, \\
		&\; \epsilon_k \in [0, \epsilon_0], \quad \forall k \in K, \\
		&\; \|\F_k\|_F^2 \le P_k, \quad \forall k \in K.
	\end{split}
\label{eq_short_blocklength_latency_min_QoS_prob}
\end{equation*}
\end{comment}

\subsection{Strategies to solve the Problem Formulation}
%******************************************************************
Solution for (\ref{eq_short_blocklength_latency_min_prob}) can be attained via alternating optimization.  First, by fixing $\F_k$, $n_k$ can be found via exhaustive search by searching the interval $\left[ n_{\min}, n_{\max} \right]$ for $k \in \mcK$.  Then after a suboptimal (because $n_k$, $\forall k$ is found based on a suboptimal $\F_k$), then fixing $\n_k$, we should find $\F_k$, for $k \in \mcK$, by solving 
\begin{equation}
	\begin{split}
		\max_{\F_k, k \in \mcK} &\; R_k(\F_k) - \sqrt{ \frac{ V_k(\F_k) }{ n_k }  } Q^{-1}(\epsilon_k), \quad \forall k \in \mcK \\
		& \; \left\| \F_k \right\|_F^2 \leq P_k, \quad \forall k \in \mcK.
	\end{split}
\label{eq_latency_min_precoder_subprob}
\end{equation}
This implies if the problem needs to be convex, then $R_k(\F_k)$ needs to be concave (where the maximization of $R_k(\F_k)$ can be done by reformulating the problem into minimization of the WMMSE) and $\sqrt{ \frac{ V_k(\F_k) }{ n_k } }$ needs to be convex.  Since square root is a monotonically nondecreasing concave function, the square root can be removed and $V_k(\F_k)$ can be approximated, e.g. using its first-order Taylor series expansion, so that (\ref{eq_latency_min_precoder_subprob}) can be convexified as
\begin{equation}
	\begin{split}
		\max_{\F_k, k \in \mcK} &\; W_k(\F_k) - \left( V_k \left( \F_k^{(p)} \right) + \nabla V_k \left( \F_k^{(p)} \right)^H  \left( \F_k - \F_k^{(p)} \right) \right), \quad \forall k \in \mcK \\
		& \; \left\| \F_k \right\|_F^2 \leq P_k, \quad \forall k \in \mcK,
	\end{split}
\label{eq_latency_min_precoder_subprob_cvx}
\end{equation}
where $\F_k^{(p)}$ is the precoder at the $p$th iteration of the successive convex optimization (SCA) (or the WMMSE update since the SCA needs to be applied in conjunction with the WMMSE iteration).  

Since the rate of the $k$th user is defined as 
\begin{equation*}
	R_k = \log \det \left(  \I_{N_R}  +  \H_k \F_k  \F_k^H \H_k^H \left(  \sigma_w^2 \I_{N_R}  + \sum_{j \neq k}  \H_j \F_j \F_j^H \H_j^H  \right)^{-1} \right), 
\end{equation*}
let $\H_k \in \complex^{N_r \times N_k}$ with $N_k$ denoting the number of transmit antennas at user $k$.  This implies $\F_k \in \complex^{N_k \times d_k}$, where $d_k \leq N_k$ denotes the number of spatial streams from the $k$th user.  Then it has been shown in \cite{qi2011} that the weighted sum rate
\begin{equation*}
	\sum_{k=1}^K w_k R_k, \quad s.t. \left\| \F_k \right\|_F^2 \leq P_k 
\end{equation*}
can be maximized by solving the WMMSE problem.  Define the $k$th receiver at the AP as $\U_K \in \complex^{N_R \times d_k}$.  Also define the weight matrix $\W_k \in \complex^{d_k \times d_k}$ where $\W_k \succ \0_{d_k \times d_k}$ (which will be equal to $\E_k^{-1}$).  Further define the MSE matrix 
\begin{equation*}
	\E_k \triangleq \I_{d_k} - \U_k^H \H_k \F_k - \F_k^H \H_k^H \U_k + \U_k^H \left( \sigma_w^2 \I_{N_r} + \sum_{j=1}^K \H_j \F_j \F_j^H \H_j^H \right) \U_k \in \complex^{d_k \times d_k}.
\end{equation*}
Then the per user matrix quadratic transform  identity implies
\begin{align*}
	&w_k \log\det\left( \I_{N_R}  +  \H_k \F_k  \F_k^H \H_k^H \left(  \sigma_w^2 \I_{N_R}  + \sum_{j \neq k}  \H_j \F_j \F_j^H \H_j^H  \right)^{-1} \right) \\
	&\quad = \max_{\U_k, \W_k \succ \0_{d_k \times d_k}} \; w_k \left( \log\det \W_k - tr\left(\W_k \E_k\right) \right) + \text{constant}.
\end{align*}
Then the original weighted sum rate prolbme $\displaystyle \max_{\F_k, k \in \mcK} \; w_k R_k$, subject to $\left\| \F_k \right\|_F^2 \leq P_k $ problem is equivalent to 
\begin{equation}
	\begin{split}
		\max_{\F_k, \W_k, \U_k, k \in \mcK} &\; w_k \left( \log\det \W_k - tr\left( \W_k \E_k \right) \right) \\
		s.t. &\; \left\| \F_k \right\|_F^2 \leq P_k.
	\end{split}
\label{eq_WMMSE_equiv_prob}
\end{equation}
Let 
\begin{equation*}
	\bSigma \triangleq \sigma_w^2 \I_{N_R} + \sum_{j=1}^K \H_j \F_j \F_j^H \H_j^H.
\end{equation*}
Then the following WMMSE algorithm can be used to find the solution for (\ref{eq_WMMSE_equiv_prob}).  
\begin{enumerate}[1.]
	\item {\bf MMSE receiver update}
	\begin{equation*}
		\U_k = \bSigma^{-1} \H_k \F_k.
	\end{equation*}
	%
	\item {\bf Weight matrix computation (inverse MSE)}
	\begin{equation*}
		\W_k = \E_k^{-1}
	\end{equation*}
	%
	\item {\bf Precoder computation (quadratic KKT with power constraint)} Let 
	\begin{align*}
		\C_k &\triangleq \H_k^H \U_k \W_k \U_k^H \H_k \in \complex^{N_k \times N_k} \\
		\D_k &\triangleq \H_k^H \U_k \W_k \in \complex^{N_k \times d_k}.
	\end{align*} 
	Then the precoder with per-user power constraint is computed as
	\begin{equation*}
		\F_k(\lambda_k) = \left( \w_k \C_k + \lambda_k \I_{N_k} \right)^{-1}  (\w_k \D_k ) \in \complex^{N_k \times d_k}, 
	\end{equation*}
	where $\lambda_k \geq 0$ is chosen by the bisection method to satisfy $tr\left( \F_k \F_k^H \right) = P_k$ if needed, otherwise $\lambda_k = 0$ and $tr\left( \F_k \F_k^H \right) < P_k$.  That is, this corresponds to the complementary slackness condition in the KKT conditions.
\end{enumerate}
Note that it has been reported that by initializing $\F_k$ using MRT precoder speeds up the convergence of the WMMSE algorithm.  Note also that the sum rate is maximized by maximizing the per-user rate $R_k$.  Note that we may be able to set the weight $w_k$ as the inverse of the CNR so that the rate of the worst-case (i.e. worst CNR) will be optimized the most.  Algorithm \ref{alg_latency_min} summarizes the steps. \\
%**************************************************************************
\begin{algorithm}[H]
	\SetAlgoLined
    \KwResult{$n_k, \; \F_k^\star, \forall k \;$}
    \SetKwInput{Input}{Initialization}
    \Input{Set $i  = 0$, choose $\varepsilon$ (e.g. small integer number), choose $\left\{ \F_k^{(0)}, \forall k \right\}$ as the MRT precoder, $w_k$ = inverse of the $\text{CNR}_k$} %\\
	\Do{ $\displaystyle \left| \sum_{k \in \mcK} n_k^{(i)} - \sum_{k \in \mcK} n_k^{(i-1)} \right| \leq \varepsilon$ }{	
		Exhaustively find $n_k$, $\forall k$ using $\F_k^{(i)}$ \;
		Denote (sub-)optimal blocklength as $n_k^\star$\;
		$n_k^{(i)} = n_k^\star$, $\forall k$ \;
		\For{ each user $k = 1, \ldots, K$ } {
			Use the WMMSE algorithm and convex approximation of $V_k(\F_k)$ to find the precoder $\F_k$ using $n_k^{(i)}$\; 
			Denote (sub-)optimal precoder as $\F_k^\star$\;
    	}
    	$i = i + 1$ \;
    	$\F_k^{(i)} = \F_k^\star$, $\forall k$ \;
    }	
    %$\hbs_W[n|n]=\hbs_W^{(i)}[n|n]$\;
    %$\M[n|n]=\left(\I-\K^{(i-1)}[n]\H \right)\M^{(i-1)}[n|n-1]$\;
    \caption{Proposed latency minimization algorithm in the uplink via precoder.}
\label{alg_latency_min}
\end{algorithm}
%**************************************************************************

%\begin{comment}
%\textbf{Problem Formulation}
%\begin{equation*}
%	\begin{split}
%		\min_{\F_k, n_k, \gamma_k, \epsilon_k} &\; \sum_{k \in K} W_k \frac{n_k}{f_s} \\
%		\text{s.t. } &\; \frac{B_k}{n_k} \le   R_k(\F_k) - \sqrt{\frac{V_k(\F_k)}{n_k}} \, Q^{-1}(\epsilon_k), \quad \forall k \in K, \\
%		&\; n_k \in [n_{\min}, n_{\max}], \quad \forall k \in K, \\
%		&\; \epsilon_k \in [0, \epsilon_0], \quad \forall k \in K, \\
%		&\; \|\F_k\|_F^2 \le P_k, \quad \forall k \in K.
%	\end{split}
%\label{eq_prob_form_basic}
%\end{equation*}
%
%\textbf{Initial Approach (8-21-8-2025)}
%\end{comment}



\subsection{Capacity and Channel Dispersion}
%****************************************************************** 
The Capacity and Channel Variance of Coherant Rayleigh Fading Channel is discussed in \cite{collins2019}. 

However \cite{collins2019} makes the folowing assumptions are made,
 
\begin{enumerate}
\item The channel is isotrophic.
\item The noise has unit variance.
\item Considers SU-MIMO.
\end{enumerate}


The Capacity of the channel of user-$k$ is defined as, 

\begin{equation*}
\begin{gathered}C(P)=\frac{1}{2}\mathbb{E}\left[\log\det\left(I_{n_{r}}+\frac{P}{n_{t}}HH^{T}\right)\right]\\=\sum_{i=1}^{n_{\mathrm{min}}}\mathbb{E}\left[C_{AWGN}\left(\frac{P}{n_{t}}\lambda_{i}^{2}\right)\right],\end{gathered}
\end{equation*}

where $C_{AWGN}(P) = \frac{1}{2}log(1+P)$ is the capacity of the AWGN with $P$ being the power and SNR (due to unit noise variance). $n_{min} = min(n_r, n_t)$, $\lambda_i^2$ $ i\in \{1, \cdots, n_{min}\}$ are eigenvalues of $\H \H^H$.

\newpage
\textbf{Channel Variance} 

Assume that $\mathbb{P}[\text{rank} \H > 1] > 0$, then $V (P) = V_{iid}(P)$, 

\begin{equation*}\begin{aligned}V_{iid}(P)&=T\mathrm{Var}\left(\sum_{i=1}^{n_{\min}}C_{AWGN}\left(\frac{P}{n_{t}}\lambda_{i}^{2}\right)\right)\\&+\sum_{i=1}^{n_{\min}}\mathbb{E}\left[V_{AWGN}\left(\frac{P}{n_{t}}\lambda_{i}^{2}\right)\right]\\&+\left(\frac{P}{n_t}\right)^2\left(\eta_1-\frac{\eta_2}{n_t}\right)\end{aligned}\end{equation*}


$\lambda_i^2$ $ i\in \{1, \cdots, n_{min}\}$ are eigenvalues of $\H \H^H$.
\begin{equation*}
V_{AWGN}(P) = \frac{log^2e}{2}(1 - \frac{1}{(1+P)^2})
\end{equation*}

\begin{equation*}
\begin{aligned}&c(\sigma)\triangleq\frac{\sigma}{1+\frac{P}{n_{t}}\sigma}\\&\eta_{1}\triangleq\frac{\log^{2}e}{2}\sum_{i=1}^{n_{\min}}\mathbb{E}\left[c^{2}(\lambda_{i}^{2})\right]\\&\eta_{2}\triangleq\frac{\log^{2}e}{2}\left(\sum_{i=1}^{n_{\min}}\mathbb{E}\left[c(\lambda_{i}^{2})\right]\right)^{2}\end{aligned}
\end{equation*}

\textbf{Obtaining a closed form}

Under high SNR $P \to \infty$. Let $n_{min} = min(n_t,n_r)$, $\alpha = P/n_t$ and $\lambda_i$ the eigenvalues of $\H \H^H$.

If we assume $P$ is a large constant ($P \to \infty$), 

\begin{equation*}
\operatorname{Var}\!\left[\sum_{i=1}^{n_{\min}} \ln\!\left(1+\alpha \lambda_i\right)\right]
\;\sim\;
\operatorname{Var}\!\left[\sum_{i=1}^{n_{\min}} \ln \lambda_i \right]
\end{equation*}

\begin{equation*}
V_{\mathrm{AWGN}}(\alpha \lambda_i) 
= (\ln e)^{2} \, \frac{(1+\alpha \lambda_i)^{2}}{\alpha \lambda_i \, (\alpha \lambda_i + 2)}
\;\sim\; (\ln e)^{2}, \quad \text{as } \alpha \to \infty
\end{equation*}

\begin{equation*}
\operatorname{Var}\!\left[\sum_{i=1}^{n_{\min}} C_{\mathrm{AWGN}}(\alpha \lambda_i)\right]
\;\sim\;
\operatorname{Var}\!\left[\sum_{i=1}^{n_{\min}} \ln(\alpha \lambda_i)\right]
= \operatorname{Var}\!\left[n_{\min}\ln \alpha + \sum_{i=1}^{n_{\min}} \ln \lambda_i\right]
= \operatorname{Var}\!\left[\sum_{i=1}^{n_{\min}} \ln \lambda_i\right]
\end{equation*}

\begin{equation*}
c(\sigma) \triangleq \frac{\sigma}{1+\frac{P}{n_t}\sigma}, 
\quad \alpha \triangleq \frac{P}{n_t}.
\end{equation*}

\textbf{Step 1: High-SNR expansion of $c(\sigma)$}

\begin{equation*}
c(\sigma) = \frac{\sigma}{1+\alpha\sigma} 
= \frac{1}{\alpha}\cdot \frac{1}{1+\tfrac{1}{\alpha\sigma}}.
\end{equation*}

Using the series $\tfrac{1}{1+x} = 1 - x + x^2 + \mathcal{O}(x^3)$ with $x=\tfrac{1}{\alpha\sigma}$:

\begin{equation*}
c(\sigma) = \frac{1}{\alpha} - \frac{1}{\alpha^2\sigma} 
+ \frac{1}{\alpha^3\sigma^2} + \mathcal{O}\!\left(\frac{1}{\alpha^4}\right).
\end{equation*}

---

\textbf{Step 2: Expansion of $\eta_1$}

\begin{equation*}
c(\sigma)^2 = \frac{1}{\alpha^2} - \frac{2}{\alpha^3\sigma} 
+ \mathcal{O}\!\left(\frac{1}{\alpha^4}\right).
\end{equation*}

\begin{equation*}
\eta_1 = \frac{\log^2 e}{2}\sum_{i=1}^{n_{\min}}
\mathbb{E}[c(\lambda_i^2)^2]
= \frac{\log^2 e}{2}\left(
\frac{n_{\min}}{\alpha^2} 
- \frac{2}{\alpha^3}\sum_{i=1}^{n_{\min}}\mathbb{E}\!\left[\tfrac{1}{\lambda_i^2}\right]
+ \mathcal{O}\!\left(\frac{1}{\alpha^4}\right)\right).
\end{equation*}

---

\textbf{Step 3: Expansion of $\eta_2$}

\begin{equation*}
\mathbb{E}[c(\lambda_i^2)] = \frac{1}{\alpha} 
- \frac{1}{\alpha^2}\mathbb{E}\!\left[\tfrac{1}{\lambda_i^2}\right]
+ \mathcal{O}\!\left(\frac{1}{\alpha^3}\right).
\end{equation*}

\begin{equation*}
\sum_{i=1}^{n_{\min}}\mathbb{E}[c(\lambda_i^2)]
= \frac{n_{\min}}{\alpha} 
- \frac{1}{\alpha^2}\sum_{i=1}^{n_{\min}}\mathbb{E}\!\left[\tfrac{1}{\lambda_i^2}\right]
+ \mathcal{O}\!\left(\frac{1}{\alpha^3}\right).
\end{equation*}

\begin{equation*}
\Bigg(\sum_{i=1}^{n_{\min}}\mathbb{E}[c(\lambda_i^2)]\Bigg)^2
= \frac{n_{\min}^2}{\alpha^2} 
- \frac{2n_{\min}}{\alpha^3}\sum_{i=1}^{n_{\min}}\mathbb{E}\!\left[\tfrac{1}{\lambda_i^2}\right]
+ \mathcal{O}\!\left(\frac{1}{\alpha^4}\right).
\end{equation*}

\begin{equation*}
\eta_2 = \frac{\log^2 e}{2}\left(
\frac{n_{\min}^2}{\alpha^2} 
- \frac{2n_{\min}}{\alpha^3}\sum_{i=1}^{n_{\min}}\mathbb{E}\!\left[\tfrac{1}{\lambda_i^2}\right]
+ \mathcal{O}\!\left(\frac{1}{\alpha^4}\right)\right).
\end{equation*}

---

\textbf{Step 4: Finding differenve of $\eta_1$ and $\eta_2$ }

\begin{equation*}
\eta_1 - \eta_2 = \frac{\log^2 e}{2}\left(
\frac{n_{\min}-n_{\min}^2}{\alpha^2}
+ \frac{2(n_{\min}-1)}{\alpha^3}\sum_{i=1}^{n_{\min}}\mathbb{E}\!\left[\tfrac{1}{\lambda_i^2}\right]
+ \mathcal{O}\!\left(\frac{1}{\alpha^4}\right)\right).
\end{equation*}

\begin{equation*}
\frac{P}{n_t}(\eta_1-\eta_2) = \alpha(\eta_1-\eta_2)
= \frac{\log^2 e}{2}\left(
\frac{n_{\min}-n_{\min}^2}{\alpha} 
+ \frac{2(n_{\min}-1)}{\alpha^2}\sum_{i=1}^{n_{\min}}\mathbb{E}\!\left[\tfrac{1}{\lambda_i^2}\right]
+ \mathcal{O}\!\left(\frac{1}{\alpha^2}\right)\right).
\end{equation*}

If we assume that $\alpha << \alpha^2 \approx \infty$,
\begin{equation*}
 \frac{P}{n_t}(\eta_1-\eta_2)  = \frac{\log^2 e}{2}\left(
\frac{n_{\min}-n_{\min}^2}{\alpha} 
+ \mathcal{O}\!\left(\frac{1}{\alpha^2}\right)\right)
\end{equation*}
---

\newpage 
\textbf{Step 5: High-SNR conclusion}

\begin{equation*}
 \frac{P}{n_t}(\eta_1-\eta_2)  = \frac{\log^2 e}{2}\left(
\frac{n_{\min}-n_{\min}^2}{\alpha} 
+ \mathcal{O}\!\left(\frac{1}{\alpha^2}\right)\right)  \;\;\longrightarrow\; 0,
\quad \text{as } P \to \infty.
\end{equation*}



Therefore,
\begin{equation*}
V(P) = V_{\mathrm{iid}}(P) \sim 
T\,\mathrm{Var}\!\left[\sum_{i=1}^{n_{\min}}\ln \lambda_i\right]
+ n_{\min}(\ln e)^2
+ \frac{\log^2 e}{2}\left(\frac{n_{\min}-n_{\min}^2}{\alpha}\right)
+ \mathcal{O}\!\left(\frac{1}{\alpha^2}\right), 
\quad \alpha=\frac{P}{n_t}.
\end{equation*}

We can further simplify this by truly considering $\alpha \to \infty$, then the last terms become $0$. However i have kept these terms in the equation to give role to the precoder $\F$ in our alternating optimization proccess.

We can further expand $T\,\mathrm{Var}\!\left[\sum_{i=1}^{n_{\min}}\ln \lambda_i^2\right]$ with the trigamma function $\psi1(x)$

\begin{equation*}
\begin{aligned}
\mathrm{Var}\!\left[\sum_{i=1}^{n_{\min}}\ln \lambda_i^2\right] 
&= \sum_{i=1}^{n_{\min}} \psi_1(n_{\max}-i+1),\\
&\quad n_{\min} = \min(n_t,n_r), \quad n_{\max} = \max(n_t,n_r),\\
&\quad \alpha = \frac{P}{n_t}
\end{aligned}
\end{equation*}

where $\psi_1(x)$ is, 
\begin{equation*}
\psi_1(x)=\sum_{k=0}^\infty\frac{1}{(x+k)^2}.
\end{equation*}

In the assumtption of very large $x$, 
\begin{equation*}
\psi_1(x)=\frac{1}{x}+\frac{1}{2x^2}+\frac{1}{6x^3}-\frac{1}{30x^5}+\frac{1}{42x^7}-\cdots
\end{equation*}

Thus when we consider $\sum_{i=1}^{n_{\min}} \psi_1(n_{\max}-i+1)$ to obtain a closed form for the trigamma function, we need to assume a very large $n_{\max}-i+1$, thus if $n_{\max} = \max(n_r,n_t) \to \infty$, 

\begin{equation*}
\psi_1(n_{\max}-i+1)=\frac{1}{n_{\max}-i+1}+\frac{1}{2(n_{\max}-i+1)^2}+\frac{1}{6(n_{\max}-i+1)^3} - \frac{1}{30(n_{\max}-i+1)^5 +} + \cdots
\end{equation*}
or using the Bernoulli numbers $B$,
\begin{equation*}
\psi_1(n_{\max}-i+1)=\sum_{k=0}^\infty\frac{B_{2k}}{(n_{\max}-i+1)^{2k+1}}
\end{equation*}

Thus we can reduce $\mathrm{Var}\!\left[\sum_{i=1}^{n_{\min}}\ln \lambda_i^2\right] $ to, 

\begin{equation*}
\mathrm{Var}\left[\sum_{i=1}^{n_{\min}}\ln\lambda_i^2\right]\approx\sum_{i=1}^{n_{\min}}\psi_1(n_{\max}-i+1)=\sum_{i=1}^{n_{\min}}\left[\frac{1}{n_{\max}-i+1}+\frac{1}{2(n_{\max}-i+1)^2}\right.+ \cdots
\end{equation*}

\newpage
\textbf{Final closed form for Channel Variance}

Thus the final $V(P)$ under the assumtion of, $P \to \infty$ and $n_{\max} \to \infty$ is, 

 \begin{equation*}
\begin{aligned}
V(P) \;=\; V_{\mathrm{iid}}(P) 
&\;\approx\; \; \sum_{i=1}^{n_{\min}}\!\left[
    \frac{1}{\,n_{\max}-i+1\,}
    + \frac{1}{2(n_{\max}-i+1)^2}
    + \cdots
\right] \\[0.6em]
&\quad +\; n_{\min}(\ln e)^2 \\[0.6em]
&\quad +\; \frac{\log^2 e}{2}\left(\frac{n_{\min}-n_{\min}^2}{\alpha}\right) \\[0.6em]
&\quad +\; \mathcal{O}\!\left(\tfrac{1}{\alpha^2}\right),
\qquad \alpha=\tfrac{P}{n_t}.
\end{aligned}
\end{equation*}


\begin{thebibliography}{25} 
\bibitem{rimoldi1996}
B. Rimoldi and R. Urbanke, ``A rate-splitting approach to the Gaussian multiple-access channel,'' \emph{IEEE Trans. Inf. Theory},  vol. 42, no. 2, pp. 364–375, Mar. 1996.

\bibitem{yang2019}
Z. Yang, M. Chen, W. Saad, W. Xu and M. Shikh-Bahaei, ``Sum-rate maximization of uplink rate splitting multiple access (RSMA) communication,'' \emph{Proc. of the IEEE Global Communications Conference (GLOBECOM)}, Waikoloa, HI, USA, pp. 1-6, 2019.  doi: 10.1109/GLOBECOM38437.2019.9013344.

\bibitem{durisi2016}
G. Durisi, T. Koch and P. Popovski, ``Toward massive, ultrareliable, and low-latency wireless communication with short packets,'' \emph{Proceedings of the IEEE}, vol. 104(9), pp. 1711-1726, 2016.

\bibitem{you2023}
X. You \emph{et al.}, ``Closed-form approximation for performance bound of finite blocklength massive MIMO transmission,'' \emph{IEEE Trans. on Communications}, vol. 71, no. 12, pp. 6939-6951, Dec. 2023

\bibitem{zhang2023}
Y. Zhang, W. Cheng, J. Wang and W. Zhang, ``Performance analysis and blocklength minimization of uplink RSMA for short packet transmissions in URLLC,'' \emph{Proc. of the IEEE Global Communications Conference (GLOBECOM)}, Kuala Lumpur, Malaysia, pp. 6765-6770, 2023.

\bibitem{qi2011}
Q. Shi, M. Razaviyayn, Z.-Q. Luo, and C. He, ``An iteratively weighted MMSE approach to distributed sum-utility maximization for a MIMO interfering broadcast channel,'' \emph{IEEE Trans. on Signal Processing}, vol. 59(9), pp. 4331–4340, Sep. 2011.

\bibitem{collins2019}
A. Collins and Y. Polyanskiy, "Coherent Multiple-Antenna Block-Fading Channels at Finite Blocklength," in IEEE Transactions on Information Theory, vol. 65, no. 1, pp. 380-405, Jan. 2019

%\bibitem{shen2018}
%K. Shen and W. Yu, ``Fractional programming for communication systems - Part I: power control and beamforming,'' \emph{IEEE Trans. on Signal Processing}, vol. 66(10), pp. 2616-2630, May 2018.

\end{thebibliography}
\end{document}
